\section{Second-Order Expansion and Risk Terms}

\subsection{IS Equation}

The nonlinear household Euler equation can be written as
\begin{equation}
    1
    =
    \E_t
    \left[
        \beta
        \frac{U_c(C_{t+1})}{U_c(C_t)}
        \frac{R_t}{\Pi_{t+1}}
    \right].
    \label{eq:nonlinear_euler}
\end{equation}
Let
\begin{equation}
    m_{t+1}
    \equiv
    \log\beta
    +
    \log U_c(C_{t+1})
    -
    \log U_c(C_t)
    +
    i_t
    -
    \pi_{t+1}.
    \label{eq:log_sdf_return}
\end{equation}
Then Equation \eqref{eq:nonlinear_euler} is
\begin{equation}
    \E_t[\exp(m_{t+1})]=1.
\end{equation}
A second-order log-normal approximation gives
\begin{equation}
    0
    \approx
    \E_t m_{t+1}
    +
    \frac{1}{2}\Var_t(m_{t+1}).
    \label{eq:second_order_euler}
\end{equation}
With CRRA utility and local notation, this implies an intertemporal demand
condition of the form
\begin{equation}
    x_t
    =
    \E_t x_{t+1}
    -
    \frac{1}{\gamma}
    \left(
        i_t-\E_t\pi_{t+1}-r^{n,1}_t
    \right)
    +
    \omega^x_t,
    \label{eq:is_with_wedge}
\end{equation}
where the demand wedge collects second-order risk corrections:
\begin{equation}
    \omega^x_t
    \equiv
    -\frac{1}{2\gamma}
    \left[
        \Var_t(m_{t+1})
        -
        \Var^n_t(m^n_{t+1})
    \right]
    +
    \Delta r^n_{t,\mathrm{risk}}.
    \label{eq:demand_wedge_definition}
\end{equation}
Here \(m^n_{t+1}\) is the flexible-price counterpart and
\(\Delta r^n_{t,\mathrm{risk}}\) denotes the correction needed when the natural
real rate is defined as a risk-adjusted object rather than as the first-order
object \(r^{n,1}_t\).

The sign in Equation \eqref{eq:demand_wedge_definition} reflects precautionary
saving: for a given real interest rate path, higher conditional marginal-utility
risk lowers current demand relative to expected future demand.

\subsection{Phillips Curve}

The New Keynesian Phillips curve is generated by a nonlinear price-setting
condition. Write that condition generically as
\begin{equation}
    F^\pi
    \left(
        \pi_t,
        x_t,
        \E_t \pi_{t+1},
        \E_t x_{t+1},
        a_t,
        \Omega_t
    \right)
    =
    0,
    \label{eq:generic_price_setting}
\end{equation}
where \(\Omega_t\) is the conditional covariance matrix of future states. A
second-order expansion around the deterministic steady state yields the usual
linear NKPC plus a curvature correction:
\begin{equation}
    \pi_t
    =
    \beta \E_t \pi_{t+1}
    +
    \kappa x_t
    +
    \omega^\pi_t.
    \label{eq:nkpc_with_wedge}
\end{equation}
The markup or price-setting wedge can be represented as
\begin{equation}
    \omega^\pi_t
    =
    \Lambda_\pi
    \left[
        \frac{1}{2}
        \tr
        \left(
            H^\pi \Omega_t
        \right)
    \right],
    \label{eq:price_wedge_definition}
\end{equation}
where \(H^\pi\) is the Hessian of the underlying price-setting condition with
respect to future states, and \(\Lambda_\pi\) converts the curvature term into
inflation units.

Equation \eqref{eq:price_wedge_definition} is intentionally generic. Without
choosing Calvo versus Rotemberg pricing, the exact sign and magnitude of
\(\omega^\pi_t\) cannot be pinned down. The important accounting point is that
TFP uncertainty can appear in the NKPC only through the curvature of price
setting, real marginal cost, and expectations.

\subsection{Policy Rule}

If the monetary policy rule is linear and contains no policy uncertainty, a TFP
volatility shock does not create a primitive policy wedge:
\begin{equation}
    \omega^i_t = 0.
    \label{eq:no_policy_wedge}
\end{equation}
For empirical accounting, however, it is useful to allow the rule to be written
as
\begin{equation}
    i_t
    =
    \rho_i i_{t-1}
    +
    (1-\rho_i)
    \left(
        \phi_\pi \pi_t+\phi_x x_t
    \right)
    +
    \omega^i_t.
    \label{eq:taylor_with_wedge}
\end{equation}
Under the single-source TFP uncertainty model, \(\omega^i_t\) should be zero
unless it is interpreted as a residual that absorbs rule misspecification,
policy smoothing not captured by the simple rule, or measurement error.
